% ========================= SECTION 3 ================================= %

\section{Agent-Based Model Equations}
\label{sec:abm_equations}

The crop--soil system is modeled as a set of $N = 130$ coupled discrete-time dynamical systems, following the ABM framework introduced by L\'{o}pez-Jim\'{e}nez et al.~\cite{Lopezjimenez2024}. Each agent $n \in \{1, \ldots, N\}$ maintains four state variables that evolve at a daily time step $k$. These states capture the essential dynamics of the soil--plant--atmosphere continuum: water availability, thermal accumulation, crop maturation, and biomass production.

\section{Agent-Based Model Equations}
\label{sec:abm_equations}

The crop--soil system is modeled as a set of $N = 130$ coupled discrete-time dynamical systems, following the ABM framework introduced by L\'{o}pez-Jim\'{e}nez et al.~\cite{Lopezjimenez2024}. Each agent $n \in \{1, \ldots, N\}$ maintains four state variables that evolve at a daily time step $k$. These states capture the essential dynamics of the soil--plant--atmosphere continuum: water availability, thermal accumulation, crop maturation, and biomass production.

\subsection{State Equations}
\label{sec:state_equations}

The four state equations governing the evolution of agent $n$ are:
\begin{align}
    x_1^n(k+1) &= x_1^n(k) - \varphi_1^n - \varphi_2^n - \varphi_3^n - \varphi_{4,\text{out}}^n + \varphi_{2,\text{in}}^n + \varphi_{4,\text{in}}^n + P(k) + u^n(k) \label{eq:x1} \\
    x_2^n(k+1) &= x_2^n(k) + h_1^n(k) \label{eq:x2} \\
    x_3^n(k+1) &= x_3^n(k) + \theta_{11}\bigl(1 - h_2(k)\bigr) + \theta_{12}\bigl(1 - h_3^n(k)\bigr) \label{eq:x3} \\
    x_4^n(k+1) &= x_4^n(k) + \theta_{13} \cdot h_3^n(k) \cdot h_6^n(k) \cdot h_7(k) \cdot g(x_2^n) \cdot R_s(k) \label{eq:x4}
\end{align}

where $x_1$ is the volumetric soil water content in the root zone (mm), $x_2$ is the cumulative thermal time ($^{\circ}$C$\cdot$day), $x_3$ is the cumulative maturation index incorporating heat and drought stress effects, and $x_4$ is the aboveground dry biomass (g/m$^2$). The environmental inputs are daily precipitation $P(k)$ (mm/day) and solar radiation $R_s(k)$ (MJ/m$^2$/day). The control input $u^n(k)$ is the irrigation applied to agent $n$ on day $k$ (mm/day).

Note that Equations \ref{eq:x1} and \ref{eq:x4} feature updated formulations compared to the baseline model to strictly enforce mass conservation and drought-yield coupling. These corrections are detailed fully in Section~\ref{sec:corrections}.

\subsection{Soil Water Fluxes}
\label{sec:soil_fluxes}

\textbf{Crop transpiration} ($\varphi_1$) represents the water extracted from the soil by plant roots to meet evaporative demand. It is modeled as a demand-limited process bounded by the reference evapotranspiration ($\text{ET}_0$):
\begin{equation}
    \varphi_1^n = \min\!\Big(\max\!\big(\theta_1(x_1^n - \theta_2 \theta_5), 0\big),\; \text{ET}_0\Big)
    \label{eq:phi1}
\end{equation}
where $\theta_1$ is the water uptake coefficient, $\theta_2$ is the volumetric wilting point, and $\theta_5$ is the root zone depth (mm). 

\textbf{Surface runoff} ($\varphi_2$) is computed using the SCS curve number method, partitioning rainfall between infiltration and overland flow when precipitation $P$ exceeds the initial abstraction threshold $\theta_3$:
\begin{equation}
    \varphi_2^n = 
    \begin{cases}
        \dfrac{(P - \theta_3)^2}{P + 4\theta_3}, & P > \theta_3 \\[8pt]
        0, & P \leq \theta_3
    \end{cases}
    \label{eq:phi2}
\end{equation}

\textbf{Deep drainage} ($\varphi_3$) represents water percolating below the root zone, activated when soil water exceeds field capacity ($\theta_6 \theta_5$):
\begin{equation}
    \varphi_3^n = \theta_4 \cdot \max\!\big(x_1^n + P - \varphi_2^n - \theta_6 \theta_5,\; 0\big)
    \label{eq:phi3}
\end{equation}
where $\theta_4$ is the drainage coefficient governing the rate of percolation.

\textbf{Lateral water exchange} ($\varphi_{4,\text{in}}$, $\varphi_{4,\text{out}}$) and \textbf{Runoff Redistribution} ($\varphi_{2,\text{in}}$) represent the topography-driven redistribution of excess soil water and surface runoff between adjacent agents based on the directed graph established in Section~\ref{sec:directed_graph}. Their formulations are detailed in Section~\ref{sec:corrections}.

\subsection{Stress and Growth Functions}
\label{sec:stress_growth}

\textbf{Thermal time accumulation} ($h_1$) drives phenological development. Only temperatures above the crop's base temperature $\theta_7$ contribute:
\begin{equation}
    h_1^n = \max\!\bigl(T_{\text{mean}} - \theta_7,\; 0\bigr)
    \label{eq:h1}
\end{equation}

\textbf{Heat stress} ($h_2$) reduces maturation efficiency linearly between an optimal temperature ceiling ($\theta_9$) and an extreme heat threshold ($\theta_{10}$), while \textbf{Low temperature stress} ($h_7$) acts as a binary switch, preventing biomass accumulation when $T_{\text{mean}} \leq \theta_7$.

\textbf{Drought stress} ($h_3$) quantifies the reduction in crop performance when actual transpiration falls short of atmospheric demand:
\begin{equation}
    h_3^n = 1 - \theta_{14} \cdot \max\!\left(1 - \frac{\varphi_1^n}{\text{ET}_0},\; 0\right)
    \label{eq:h3}
\end{equation}
where $\theta_{14}$ controls the crop's sensitivity to drought. 

\textbf{Waterlogging stress} ($h_6$) penalizes biomass growth when soil water exceeds field capacity, simulating oxygen deprivation in the root zone:
\begin{equation}
    h_6^n = \text{clip}\!\left(1 - \frac{x_1^n - \theta_6 \theta_5}{\theta_6 \theta_5},\; 0,\; 1\right)
    \label{eq:h6}
\end{equation}

\textbf{Growth function} ($g$) determines the fraction of intercepted radiation contributing to biomass, following a two-branch sigmoid model driven by cumulative thermal time:
\begin{equation}
    g(x_2^n) = 
    \begin{cases}
        \dfrac{\theta_{19}}{1 + \exp\!\bigl(-0.01(x_2^n - \theta_{20})\bigr)}, & x_2^n \leq \dfrac{\theta_{18}}{2} \\[14pt]
        \dfrac{\theta_{19}}{1 + \exp\!\bigl(0.01(x_2^n + x_3^n - \theta_{18})\bigr)}, & x_2^n > \dfrac{\theta_{18}}{2}
    \end{cases}
    \label{eq:growth}
\end{equation}
The first branch governs the vegetative phase, expanding sigmoidally as leaf area increases. The second branch governs senescence, declining as the cumulative stress state ($x_3$) pushes the crop toward maturity ($\theta_{18}$).

\section{Model Corrections and Extensions}

Three corrections to the baseline \cite{Lopezjimenez2024} were implemented.

\subsection{Mass Conservation}

The original water balance includes $\varphi_{4,\text{in}}$ but omits $\varphi_{4,\text{out}}$, violating mass conservation. The corrected Eq.~\ref{eq:water} includes $-\varphi_{4,\text{out}}^n$.

\subsection{Drought Stress on Biomass}

The original biomass equation lacks $h_3$, so drought does not reduce biomass. The corrected Eq.~\ref{eq:biomass} includes $h_3^n$, reducing biomass by up to $\theta_{14} \times 100\%$ under maximum drought (80\% for rice, 60\% for tobacco). This is essential for irrigation optimization: without it, the controller has no incentive to irrigate.

\subsection{Surface Runoff Redistribution}
\label{sec:runoff}

Originally, $\varphi_2$ leaves the system. A simple redistribution was implemented: each agent's runoff flows to lower neighbors with immediate infiltration ($\varphi_{2,\text{in}}^n$). A cascading alternative (top-to-bottom SCS re-application) was also tested.

\begin{table}[H]
\centering
\caption{Runoff redistribution comparison, no irrigation, wet scenario (2024).}
\label{tab:runoff}
\begin{tabular}{lccc}
\toprule
\textbf{Mode} & \textbf{Yield (kg/ha)} & \textbf{$x_1$ sinks (mm)} & \textbf{$x_1$ hills (mm)} \\
\midrule
None (original) & 2135 & 71.2 & 71.2 \\
Simple (A) & 2273 & 74.8 & 72.0 \\
Cascade (B) & 2278 & 79.5 & 72.0 \\
\bottomrule
\end{tabular}
\end{table}

Simple redistribution captures the first-order effect (6.5\% yield improvement) with negligible difference from cascade (0.2\%). Simple was adopted.

\begin{figure}[H]
\centering
\includegraphics[width=\textwidth]{figures/fig_runoff_comparison.png}
\caption{No-irrigation runoff mode comparison, wet scenario (2024).}
\label{fig:runoff}
\end{figure}

\subsection{Biomass Calibration}

$\theta_{13}$ was calibrated so $x_4$ represents g/m$^2$ dry matter. Yield: $Y = x_4 \times \text{HI} \times 10$ (kg/ha). Rice: validated against $\sim$2000~kg/ha (farmer reports) and 3500--4000~kg/ha (optimized)~[14]. Initial conditions: rice $x_{2,\text{init}} = 210$, $x_{4,\text{init}} = 60$; tobacco $x_{2,\text{init}} = 0$, $x_{4,\text{init}} = 8$.


\section{Model Validation}

\subsection{GWETROOT Cross-Validation}

ABM soil water was compared to NASA MERRA-2 GWETROOT. Both normalized: ABM as $(x_1 - \text{WP})/(\text{FC} - \text{WP})$; GWETROOT as native 0--1 fraction. Pearson correlation assessed temporal pattern agreement.

\begin{table}[H]
\centering
\caption{Cross-validation: ABM vs NASA GWETROOT.}
\label{tab:crossval}
\begin{tabular}{lccc}
\toprule
\textbf{Scenario} & \textbf{Year} & \textbf{Pearson $r$} & \textbf{$p$-value} \\
\midrule
Dry & 2022 & 0.738 & $3.1 \times 10^{-17}$ \\
Moderate & 2020 & 0.534 & $3.4 \times 10^{-8}$ \\
Wet & 2024 & 0.337 & $9.6 \times 10^{-4}$ \\
\bottomrule
\end{tabular}
\end{table}

\begin{figure}[H]
\centering
\includegraphics[width=\textwidth]{figures/fig_gwetroot_timeseries.png}
\caption{Normalized soil moisture: ABM (blue) vs NASA GWETROOT (brown).}
\label{fig:gwetroot}
\end{figure}

The dry-year $r = 0.74$ ($p < 0.001$) confirms the ABM captures wetting/drying dynamics. Decreasing correlation with increasing rainfall is expected (reduced signal contrast in wet conditions). The absolute level difference (ABM $\sim$0.08 vs GWETROOT $\sim$0.46) reflects: (1) NASA's deeper root zone ($\sim$1000~mm vs 400~mm); (2) GWETROOT observes irrigated fields while the ABM simulates rain-fed conditions; (3) different soil parameterizations.